% Generated draft: requires mathematical and source review.
\documentclass[11pt]{article}
\usepackage[margin=25mm]{geometry}
\usepackage{fontspec}
\setmainfont{Latin Modern Roman}
\usepackage{amsmath,amssymb,mathtools,unicode-math}
\setmathfont{Latin Modern Math}
\usepackage{enumitem,needspace,xurl,hyperref}
\setlength{\parindent}{0pt}
\setlength{\parskip}{5pt}
\setlength{\emergencystretch}{3em}
\newcommand{\N}{\mathbb{N}}
\newcommand{\Z}{\mathbb{Z}}
\newcommand{\Q}{\mathbb{Q}}
\newcommand{\R}{\mathbb{R}}
\newcommand{\C}{\mathbb{C}}
\newcommand{\F}{\mathbb{F}}
\newcommand{\A}{\mathbb{A}}
\newcommand{\PP}{\mathbb{P}}
\newcommand{\E}{\mathbb{E}}
\newcommand{\eps}{\varepsilon}
\newcommand{\dd}{\,\mathrm{d}}
\begin{document}
\section*{UnsolvedMath ID: 800661}
\textbf{Source problem code:} catalogue entry
\subsection{Definitions and mathematical statement}
Let \(X\) be a topological space. To make the question explicit, use the conventional interpretation
\[
\mathcal D(X)=\left\{\mathcal U:\mathcal U\text{ is a family of nonempty open subsets of }X,\quad \overline{\bigcup\mathcal U}=X\right\}.
\]
The supplied excerpt does not separately define \(\mathcal D\); this convention requires confirmation against the source. Define, as in the excerpt,
\[
\mathcal K(X)=\left\{\mathcal U\in\mathcal D(X):X=\bigcup_{U\in\mathcal U}\overline U\right\},
\]
where closures are taken in \(X\).

Under the usual convention for \(\mathsf G_1(\mathcal K(X),\mathcal D(X))\), in inning \(n\in\{1,2,\ldots\}\), ONE chooses \(\mathcal U_n\in\mathcal K(X)\), and TWO chooses \(U_n\in\mathcal U_n\). TWO wins exactly when \(\{U_n:n\geq1\}\in\mathcal D(X)\); otherwise ONE wins. A winning strategy for ONE prescribes legal moves from the preceding finite play and wins against every legal sequence of responses by TWO. These game conventions are not explicitly given in the supplied excerpt.

A Tychonoff space means a completely regular Hausdorff space. Its cellularity is
\[
c(X)=\sup\left\{|\mathcal V|:\mathcal V\text{ is a pairwise disjoint family of nonempty open subsets of }X\right\}.
\]
The source attributes to Scheepers the definition
\[
\mathfrak j=\min\left\{\kappa:\begin{gathered}
\kappa\text{ is a cardinal, and for every Tychonoff space }X,\\
c(X)\geq\kappa\ \Longrightarrow\ \text{ONE has a winning strategy in }
\mathsf G_1(\mathcal K(X),\mathcal D(X))
\end{gathered}\right\}.
\]

For interpreting the reported bounds, let \(\mathcal M\) be the ideal of meager subsets of \(\R\), and let
\[
\mathsf{cov}(\mathcal M)=\min\{|\mathcal A|:\mathcal A\subseteq\mathcal M,\ \bigcup\mathcal A=\R\}.
\]
Under the standard real-line interpretation, a set \(A\subseteq\R\) has strong measure zero if, for every sequence \((\varepsilon_n)_{n\geq1}\) of positive real numbers, there are intervals \(I_n\) of length less than \(\varepsilon_n\) with \(A\subseteq\bigcup_{n\geq1}I_n\). Write
\[
\mathsf{non}(SMZ)=\min\{|A|:A\subseteq\R\text{ does not have strong measure zero}\}.
\]
The excerpt does not specify these underlying conventions.

The supplied source reports that Tkachuk proved that the Continuum Hypothesis, \(2^{\aleph_0}=\aleph_1\), implies that ONE has a winning strategy in this game on any space of uncountable cellularity. No additional hypotheses on those spaces are stated in the excerpt. It also reports that Scheepers proved
\[
\mathsf{cov}(\mathcal M)\leq\mathfrak j\leq\mathsf{non}(SMZ).
\]
These are source-attributed background assertions, not independently verified results here.

The question, identified as Problem 34 and attributed there to [coc4, Problem 1], is: Is \(\mathfrak j\) equal to any standard cardinal characteristic of the continuum? The phrase ``standard cardinal characteristic'' does not specify a collection of candidate invariants, and the intended axioms for an equality are not explicitly stated.
\subsection{Short English statement}
The invariant \(\mathfrak j\) is the least cellularity threshold such that ONE can force TWO's chosen open sets to have nondense union in the stated game on every Tychonoff space meeting that threshold. The problem asks whether \(\mathfrak j\) coincides with a standard cardinal characteristic of the continuum. The supplied registry labels the question open; this status and the reported bounds have not been independently verified.
\subsection{Sources}
\begin{itemize}
\item Local UnsolvedMath registry, permanent ID \(800661\), source problem code catalogue entry, AMR Open Problem Lists. The supplied record is titled ``Selection Principles Milestone Problems --- Problem 34 ([coc4, Problem 1])'' and labels the problem open. Its proposed-by and proposed-year fields are Tsaban and 2003.
\item Tsaban, \emph{Selection Principles in Mathematics: A Milestone of Open Problems} (2003), Problem 34. Supplied source URL: \url{https://arxiv.org/abs/math/0312182}. The record reports access on 2026-10-16 and extraction from arXiv LaTeX; these are supplied provenance statements, not a claim of access or independent verification here.
\item [coc4, Problem 1]: supplied attribution and context. The excerpt attributes the definition of \(\mathfrak j\) and the displayed bounds to Scheepers. No full bibliographic entry or separate URL for [coc4] is supplied.
\item [tkachuk]: supplied reference label for the result attributed to Tkachuk concerning the Continuum Hypothesis and spaces of uncountable cellularity. No full bibliographic entry or separate URL is supplied.
\end{itemize}
\end{document}
